\centerline{\bf \Huge Стратегия выигрыша}
\title{Стратегия выигрыша}

\begin{flushright}
    \scriptsize Неисповедима и священна тайна равновесия весов твоих, о судьба!\\Виктор Гюго. Отверженные.
\end{flushright}

\problem{0.1} Имеются чашечные весы без гирь и 3 одинаковые по внешнему виду монеты, одна из которых фальшивая: она легче настоящих (настоящие монеты одного веса). Сколько надо взвешиваний, чтобы определить фальшивую монету? 

\problem{0.2} Как при помощи чашечных весов без гирь разделить 24 кг гвоздей на две части — 9 и 15 кг?

\problem{0.3} Золотоискатель Джек добыл 9 кг золотого песка. Сможет ли он за три взвешивания отмерить 2 кг песка с помощью чашечных весов: а) с двумя гирями — 200 г и 50 г; б) с одной гирей 200 г?

\problem{1} В корзине лежат 13 яблок. Имеются весы, с помощью которых можно узнать суммарный вес любых двух яблок.
Придумайте способ выяснить за 8 взвешиваний суммарный вес всех яблок.
\problem{2} Какие веса могут иметь три гири для того, чтобы с их помощью можно было взвесить любое целое число килограммов от 1 до 10 на чашечных весах (гири можно ставить на обе чашки)? Приведите пример.

\problem{3} Четырьмя гирями продавец может взвесить любое целое число килограммов, от 1 до 40 включительно. Общая масса гирь равна 40 кг. Какими гирями располагает продавец?

\problem{4} Семь монет расположены по кругу. Известно, что какие-то четыре из них, идущие подряд, – фальшивые и что каждая фальшивая монета легче настоящей. Объясните, как найти две фальшивые монеты за одно взвешивание на чашечных весах без гирь. (Все фальшивые монеты весят одинаково.)

\problem{5} На физическом кружке учитель поставил следующий эксперимент. Он разложил на чашечные весы 16 гирек массами 1, 2, 3, ..., 16 грамм так, что одна из чаш перевесила. Пятнадцать учеников по очереди выходили из класса и забирали с собой по одной гирьке, причем после выхода каждого ученика весы меняли свое положение и перевешивала противоположная чаша весов. Какая гирька могла остаться на весах?

\newpage

\problem{6} Известно, что среди нескольких монет имеется ровно одна фальшивая (отличается по весу от настоящих). С помощью двух взвешиваний на чашечных весах без гирь определите, легче или тяжелее фальшивая монета настоящей (находить ее не надо), если монет:\\
а) 100;\\
б) 99;\\
в) 98?
  
\problem{7} Дан мешок сахарного песка, чашечные весы и гирька в 1 г. Можно ли за 10 взвешиваний отмерить 1 кг сахара?